\documentclass[11pt,a4paper]{article}

\usepackage[utf8]{inputenc}
\usepackage[T1]{fontenc}
\usepackage{lmodern}
\usepackage[margin=1in]{geometry}
\usepackage{amsmath,amssymb}
\usepackage{enumitem}
\usepackage{booktabs}
\usepackage{fancyhdr}
\usepackage{xcolor}
\usepackage{mdframed}

\definecolor{solngray}{gray}{0.95}
\newmdenv[backgroundcolor=solngray,linewidth=0.5pt,linecolor=gray,innertopmargin=8pt,innerbottommargin=8pt]{solnbox}

\pagestyle{fancy}
\fancyhf{}
\fancyhead[L]{\textbf{Discrete Structures I}}
\fancyhead[R]{\textbf{Answer Key - Exam E}}
\fancyfoot[C]{\thepage}
\renewcommand{\headrulewidth}{0.5pt}

\begin{document}

\begin{center}
{\LARGE\textbf{ANSWER KEY}}\\[5pt]
{\Large Prelim Practice Exam --- Version E}\\[10pt]
\rule{0.8\textwidth}{0.5pt}
\end{center}

\vspace{15pt}

%==============================================================================
% PART I: MULTIPLE CHOICE
%==============================================================================
\section*{PART I: Multiple Choice (10 points)}

\textbf{1. Answer: (B)}

\begin{solnbox}
$p \to q \equiv \neg p \lor q$

This is the \textbf{Implication Law} --- the most important equivalence for conditionals!

\textbf{Quick verification:}
\begin{itemize}[noitemsep]
    \item When $p = T$, $q = F$: $p \to q = F$, and $\neg p \lor q = F \lor F = F$ \checkmark
    \item When $p = F$: $p \to q = T$, and $\neg p \lor q = T \lor q = T$ \checkmark
\end{itemize}

Note: (D) $\neg q \lor \neg p$ is equivalent to $q \to p$ (the converse), NOT the original.
\end{solnbox}

\vspace{10pt}

\textbf{2. Answer: (D)}

\begin{solnbox}
``Some students passed'' = ``There exists someone who is a student AND passed.''

$$\exists x\, (S(x) \land P(x))$$

\textbf{Critical rule:}
\begin{itemize}[noitemsep]
    \item ``ALL A are B'' uses IMPLICATION: $\forall x\, (A(x) \to B(x))$
    \item ``SOME A are B'' uses CONJUNCTION: $\exists x\, (A(x) \land B(x))$
\end{itemize}

Why? For ``all,'' we're saying IF something is an A, THEN it's a B. For ``some,'' we're saying there EXISTS something that IS BOTH an A AND a B.
\end{solnbox}

\vspace{10pt}

\textbf{3. Answer: (C)}

\begin{solnbox}
For $p \to q$:
\begin{itemize}[noitemsep]
    \item \textbf{Contrapositive:} $\neg q \to \neg p$ (equivalent to original)
    \item \textbf{Converse:} $q \to p$ (NOT equivalent)
    \item \textbf{Inverse:} $\neg p \to \neg q$ (NOT equivalent)
\end{itemize}

The inverse negates both parts but keeps the order.

\textbf{Memory tip:} Inverse and Converse are equivalent to each other, but neither is equivalent to the original. Only the Contrapositive is equivalent to the original.
\end{solnbox}

\vspace{10pt}

\textbf{4. Answer: (C)}

\begin{solnbox}
\begin{itemize}[noitemsep]
    \item \textbf{Tautology:} Always TRUE (e.g., $p \lor \neg p$)
    \item \textbf{Contradiction:} Always FALSE (e.g., $p \land \neg p$)
    \item \textbf{Contingency:} Sometimes true, sometimes false (depends on variable values)
    \item \textbf{Fallacy:} An invalid argument form (not a type of proposition)
\end{itemize}
\end{solnbox}

\vspace{10pt}

\textbf{5. Answer: (B)}

\begin{solnbox}
$\dfrac{p \land q}{\therefore p}$ is \textbf{Simplification}.

If a conjunction is true, each part is individually true. We ``simplify'' by extracting one part.

Compare with:
\begin{itemize}[noitemsep]
    \item \textbf{Conjunction:} $\dfrac{p \quad q}{\therefore p \land q}$ (opposite of simplification)
    \item \textbf{Addition:} $\dfrac{p}{\therefore p \lor q}$
\end{itemize}
\end{solnbox}

\newpage

%==============================================================================
% PART II: SHORT ANSWER
%==============================================================================
\section*{PART II: Short Answer (20 points)}

\textbf{Question 1.} (5 points)

\begin{solnbox}
\textbf{Truth table for $(p \land q) \to (p \lor q)$:}

\begin{center}
\begin{tabular}{cc|c|c|c}
\toprule
$p$ & $q$ & $p \land q$ & $p \lor q$ & $(p \land q) \to (p \lor q)$ \\
\midrule
T & T & T & T & T \\
T & F & F & T & T \\
F & T & F & T & T \\
F & F & F & F & T \\
\bottomrule
\end{tabular}
\end{center}

\textbf{Result:} All rows are TRUE, so this is a \textbf{TAUTOLOGY}.

\textbf{Why does this make sense?} If both $p$ and $q$ are true ($p \land q$), then certainly at least one of them is true ($p \lor q$). The conjunction is ``stronger'' than the disjunction.

Remember: An implication with a false antecedent is always true (rows 2, 3, 4 all have $p \land q = F$).
\end{solnbox}

\vspace{15pt}

\textbf{Question 2.} (5 points)

\begin{solnbox}
Breaking down the sentence:

\textbf{Part 1:} ``If John studies hard AND attends all classes, THEN he will pass''

$(s \land a) \to p$

\textbf{Part 2:} ``but if he doesn't study, he will fail''

$\neg s \to f$

\textbf{Combining with ``but'':} ``But'' here means ``and'' --- both conditions hold.

$$\boxed{((s \land a) \to p) \land (\neg s \to f)}$$

\textbf{Note:} Some might argue ``fail'' is $\neg p$, making part 2 become $\neg s \to \neg p$. Both interpretations are valid depending on whether $f$ and $\neg p$ are the same concept.
\end{solnbox}

\vspace{15pt}

\textbf{Question 3.} (4 points)

\begin{solnbox}
\textbf{(i) $\forall n\, (n \leq n^2)$ --- Domain: integers}

Is every integer $\leq$ its square?

Counterexample: $n = -1$. Is $-1 \leq (-1)^2 = 1$? Yes, $-1 \leq 1$. \checkmark

Wait, let's try $n = 0$: $0 \leq 0$? Yes. \checkmark

What about negative numbers in general? For $n < 0$: $n < 0 \leq n^2$ always. \checkmark

For $n \geq 0$: When $n = 0$ or $n \geq 1$, we have $n \leq n^2$. \checkmark

\textbf{Answer: TRUE}

\rule{\linewidth}{0.3pt}

\textbf{(ii) $\exists m \exists n\, (m + n < m)$ --- Domain: integers}

Is there a pair of integers where $m + n < m$?

Simplify: $m + n < m \Rightarrow n < 0$

So we just need $n$ to be negative. Example: $m = 5$, $n = -3$. Then $5 + (-3) = 2 < 5$. \checkmark

\textbf{Answer: TRUE}
\end{solnbox}

\vspace{15pt}

\textbf{Question 4.} (6 points)

\begin{solnbox}
\textbf{(i) Nobody has sent a letter to everyone.}

``There does not exist a person $x$ such that for all people $y$, $x$ sent a letter to $y$.''

$$\neg \exists x \forall y\, L(x, y) \equiv \boxed{\forall x \exists y\, \neg L(x, y)}$$

Simplified form: ``For every person, there's someone they haven't written to.''

\rule{\linewidth}{0.3pt}

\textbf{(ii) Everyone has sent a letter to someone.}

``For every person $x$, there exists a person $y$ such that $x$ sent a letter to $y$.''

$$\boxed{\forall x \exists y\, L(x, y)}$$

\rule{\linewidth}{0.3pt}

\textbf{(iii) There is someone who has received letters from everyone.}

``There exists a person $y$ such that for all people $x$, $x$ sent a letter to $y$.''

$$\boxed{\exists y \forall x\, L(x, y)}$$

\textbf{Note:} In $L(x, y)$, $x$ is the sender and $y$ is the recipient. So ``received from everyone'' means everyone sent TO that person.
\end{solnbox}

\newpage

%==============================================================================
% PART III: PROBLEM SOLVING
%==============================================================================
\section*{PART III: Problem Solving (25 points)}

\textbf{Question 5.} (6 points)

\begin{solnbox}
\textbf{(i)} $p$, therefore $p \lor q$

This is \textbf{Addition}!

$$\frac{p}{\therefore p \lor q}$$

If something is true, we can OR it with anything and it's still true.

\rule{\linewidth}{0.3pt}

\textbf{(ii)} $p \to q$, $p$, therefore $q$

This is \textbf{Modus Ponens}!

$$\frac{p \to q \quad p}{\therefore q}$$

The antecedent is true, so the consequent must follow.

\rule{\linewidth}{0.3pt}

\textbf{(iii)} $p \to q$, $q \to r$, therefore $p \to r$

This is \textbf{Hypothetical Syllogism}!

$$\frac{p \to q \quad q \to r}{\therefore p \to r}$$

Chain of implications: snow $\to$ icy roads $\to$ schools close.
\end{solnbox}

\vspace{15pt}

\textbf{Question 6.} (7 points)

\begin{solnbox}
Prove: $(p \to q) \lor (p \to r) \equiv p \to (q \lor r)$

\begin{align*}
(p \to q) \lor (p \to r) &\equiv (\neg p \lor q) \lor (\neg p \lor r) && \text{[Implication Law, twice]} \\
&\equiv \neg p \lor (q \lor r) \lor \neg p && \text{[Associative Law - rearranging]} \\
&\equiv \neg p \lor \neg p \lor (q \lor r) && \text{[Commutative Law]} \\
&\equiv \neg p \lor (q \lor r) && \text{[Idempotent Law: } \neg p \lor \neg p = \neg p\text{]} \\
&\equiv p \to (q \lor r) && \text{[Implication Law]}
\end{align*}

$$\boxed{(p \to q) \lor (p \to r) \equiv p \to (q \lor r)}$$

\textbf{Cleaner approach:}
\begin{align*}
(\neg p \lor q) \lor (\neg p \lor r) &\equiv \neg p \lor q \lor \neg p \lor r && \text{[Associative]} \\
&\equiv \neg p \lor \neg p \lor q \lor r && \text{[Commutative]} \\
&\equiv \neg p \lor (q \lor r) && \text{[Idempotent]} \\
&\equiv p \to (q \lor r) && \text{[Implication]}
\end{align*}

\textbf{Laws used:} Implication, Associative, Commutative, Idempotent
\end{solnbox}

\vspace{15pt}

\textbf{Question 7.} (4 points)

\begin{solnbox}
Original: $\forall x \exists y\, (P(x,y) \land Q(x,y))$

where $P(x,y)$: $x + y > 0$ and $Q(x,y)$: $xy > 0$.

\textbf{Negate:}
\begin{align*}
&\neg \forall x \exists y\, (P(x,y) \land Q(x,y)) \\
&\equiv \exists x\, \neg \exists y\, (P(x,y) \land Q(x,y)) && \text{[Flip } \forall \text{ to } \exists\text{]} \\
&\equiv \exists x \forall y\, \neg(P(x,y) \land Q(x,y)) && \text{[Flip } \exists \text{ to } \forall\text{]} \\
&\equiv \exists x \forall y\, (\neg P(x,y) \lor \neg Q(x,y)) && \text{[De Morgan's Law]}
\end{align*}

$$\boxed{\exists x \forall y\, (\neg P(x,y) \lor \neg Q(x,y))}$$

\textbf{In English:} ``There exists a real number $x$ such that for all real numbers $y$, either $x + y \leq 0$ or $xy \leq 0$.''
\end{solnbox}

\vspace{15pt}

\textbf{Question 8.} (8 points)

\begin{solnbox}
\textbf{Premises:}
\begin{enumerate}[noitemsep]
    \item $\neg p \lor q$
    \item $r \to p$
    \item $r$
\end{enumerate}

\textbf{Goal:} Prove $q$

\textbf{Strategy:} Note that $\neg p \lor q \equiv p \to q$. So we have $p \to q$ and need to get $p$ to apply Modus Ponens.

\begin{center}
\begin{tabular}{c|l|l}
\toprule
\textbf{Step} & \textbf{Statement} & \textbf{Reason} \\
\midrule
1 & $r \to p$ & Premise \\
2 & $r$ & Premise \\
3 & $p$ & Modus Ponens (1, 2) \\
4 & $\neg p \lor q$ & Premise \\
5 & $p \to q$ & Implication equivalence (4) [or keep as is] \\
6 & $q$ & Modus Ponens (3, 5) \\
\bottomrule
\end{tabular}
\end{center}

\textbf{Alternative without step 5:}

From step 3, $p$ is true, so $\neg p$ is false.

From step 4, $\neg p \lor q$ with $\neg p$ false means $q$ must be true (Disjunctive Syllogism).

$$\boxed{\therefore q}$$

\textbf{Explanation:}
\begin{itemize}[noitemsep]
    \item Steps 1-3: From $r$ and $r \to p$, we get $p$ by Modus Ponens.
    \item Steps 4-6: Recognize that $\neg p \lor q$ is equivalent to $p \to q$. With $p$ and $p \to q$, we conclude $q$.
\end{itemize}
\end{solnbox}

\vspace{20pt}

\begin{center}
\rule{0.5\textwidth}{0.5pt}\\[5pt]
\textbf{END OF ANSWER KEY}
\end{center}

\end{document}
